% =====================================================================
% CAPITOLO 3 — ARCHITETTURA DEL PORTALE E DEGLI AGENTI VERTICALI
% =====================================================================
% Parte A1 della "polpa". Scopo: decisioni progettuali di altro livello e motivazioni.
% Il "come" tecnico è nel Cap. 4.
% =====================================================================

% \chapter{} apre un nuovo capitolo numerato e aggiunge la voce al ToC.
% \label{} associa un'etichetta che altri comandi \ref{} possono citare.

\chapter{La suite di agenti: architettura}
\label{chap:agenti-architettura}

Questo capitolo si concentra sulla prima fase del progetto, descrivendo le scelte architetturali compiute, e i vincoli e requisiti che le hanno condizionate.
In particolare sono descritte ad alto livello le architetture dei singoli agenti verticali, dei sistemi che li ospitano e dei
flussi di rilascio e gestione del ciclo di vita del software;
i dettagli implementativi sono invece riportati nel Capitolo~\ref{chap:agenti-implementazione}.

% =====================================================================
\section{Principi guida e architettura d'insieme}
\label{sec:agenti-principi}
% =====================================================================

\subsection{Principi guida}
\label{subsec:agenti-principi-guida}

La progettazione è stata guidata da una serie di principi fondamentali trasversali,
stabiliti prima di avviare lo sviluppo e mantenuti durante tutte le fasi successive.

\begin{description}

  \item[Rispetto dei principi etici e cooperativi]
    Gli agenti LLM sono intrinsecamente portati a esprimere bias e alle allucinazioni: la letteratura recente mostra come il fenomeno allucinatorio sia un limite strutturale dei modelli generativi, non eliminabile per via puramente algoritmica~\cite{xu2024hallucination}.
    Per Coop Alleanza è importante mitigare i possibili effetti negativi sugli utenti, caratterizzando inoltre il comportamento degli agenti al pieno rispetto dei principi etici che sono alla base della cooperativa.
    Mettendo a disposizione degli utenti un pacchetto di assistenti per le loro necessità, si può prevenire l'utilizzo da parte loro di strumenti AI esterni pubblici e non autorizzati (fenomeno chiamato "Shadow AI").
    Oltre alla riduzione degli evidenti rischi per la sicurezza del dato, gli strumenti a sviluppo interno possono essere "istruiti" secondo dei "prompt di sistema" stabiliti collegialmente fra i vari uffici interni, facendo sì
    che le risposte prodotte siano conformi ai valori cooperativi ed etici che animano la cooperativa, mantenendo inoltre la possibilità di modifiche in corso d'opera e cicli di aggiornamento rapidi e controllati.

  \item[Governance facilitata e efficace]
    Gli agenti devono implementare soluzioni adeguate a prevenire utilizzi impropri e fughe di dati personali.
    Le configurazioni relative al prompt di sistema (la "personalità" dell'agente) e del grado di controllo su utilizzi impropri e dati personali devono essere visualizzabili e impostabili da una dashboard apposita,
    che consenta agli uffici preposti (che non hanno competenza tecnica stretta) di controllare questi parametri. Inoltre deve essere possibile, nel rispetto della privacy, monitorare i casi di eventi segnalati come "pericolosi".
    Tutto il codice sorgente deve essere presente e aggiornato nel repository Git usato dalla cooperativa.

  \item[Implementazione agnostica e sostenibile]
    Nella realtà aziendale attuale, molti fornitori propongono i propri modelli di assistenti o agenti AI.
    Spesso questo comporta il rischio di una dipendenza (vendor lock-in) più o meno rilevante: ad esempio, le soluzioni "out of the box" risultano di difficile trasposizione verso realtà infrastrutturali
    diverse da quelle del fornitore. Inoltre spesso sono richiesti piani di abbonamento o sottoscrizioni che possono comportare costi fissi elevati anche in caso di effettivo inutilizzo dello strumento.
    La soluzione proposta si prefigge di essere completamente portabile e facilmente adattabile verso qualunque architettura cloud o on-premises, e inoltre vuole attuare un controllo attivo dei costi, stanziando
    periodicamente un budget di utilizzo definito dagli amministratori ("token") per ogni utente, così da garantire la sostenibilità del progetto.

  \item[Sovranità del dato]
    Nessun dato deve lasciare il perimetro cloud europeo, in conformità al Regolamento Generale sulla Protezione dei Dati (GDPR).

  \item[Modularità verticale]
    Ogni dominio informativo o funzionale è affidato a un agente autonomo e indipendente,
    con il proprio meccanismo di grounding, il proprio ciclo di vita applicativo e la propria infrastruttura di deployment: in questo modo è possibile effettuare \emph{fine-tuning} o \emph{few-shot learning}~\cite{brown2020gpt3} ad hoc per il singolo caso d'uso.
    Essendo questa la modalità principale di fruizione degli agenti LLM disponibili, la puntualità e precisione delle configurazioni può impattare considerevolmente le performance.
    La separazione a livello sistemistico consente inoltre di aggiornare, monitorare e scalare ogni servizio in modo indipendente.

  \item[Automazione dei rilasci]
    Le evolutive devono potere essere rilasciate in maniera trasparente verso l'utente, automatizzata per ridurre al minimo la possibilità di errore umano, e con uno step approvativo per il passaggio in produzione.

  \item[Accesso privato e con controllo centralizzato]
    L'intero ecosistema deve potere essere accessibile solo da rete interna di Coop Alleanza, o da VPN: l'infrastruttura utilizzata deve quindi evitare l'esposizione su rete pubblica.
    Il controllo dell'abilitazione degli utenti all'uso degli assistenti deve integrarsi con il sistema di Identity Management già esistente in cooperativa, basato sui protocolli OAuth~2.0~\cite{rfc6749oauth2} e OpenID Connect~\cite{oidc1_0core} su Microsoft Entra ID~\cite{msft_entra_id}.

\end{description}

\subsection{Architettura d'insieme}
\label{subsec:agenti-architettura-insieme}

L'ecosistema di agenti deve essere accessibile da un unico "portale", così che l'utente possa abituarsi all'utilizzo di un unico indirizzo di accesso: è la logica conseguenza della modularità verticale descritta.
I nuovi agenti potranno semplicemente essere inseriti nel "framework" esistente, andando a impostare granularmente quali agenti sono "visibili" tramite i parametri dell'utenza di dominio associata all'utente.
Lo stesso approccio deve valere per le dashboard di configurazione (gli "admin panel"), anch'essi radunati in un portale dedicato.

Pur mantenendo sempre saldo il principio della maggiore agnosticità e indipendenza possibile, tutte le implementazioni sono state fatte avvalendosi delle risorse offerte dalla piattaforma Cloud di Google~\cite{gcp_vertexai,gcp_cloudrun}:
si tratta già del fornitore principale della Cooperativa in quest'ottica.


\begin{figure}[htbp]
  \centering
  \fbox{\parbox{0.88\textwidth}{\centering\vspace{2.5cm}
      \textit{[Diagramma architetturale --- da inserire]}\\[0.5em]
      \small Portal $\rightarrow$ Entra ID $\rightarrow$ Badge
      $\rightarrow$ Agenti Cloud Run\\
      (Vera, CooPolicy, AllyCare, SysAid)\\[0.3em]
      Infrastruttura condivisa: Vertex AI, BigQuery, Cloud SQL
      \vspace{2.5cm}}}
  \caption[Architettura d'insieme della suite di agenti]{Architettura
    della suite di agenti verticali di Coop Alleanza~3.0: portal con
    autenticazione federata, quattro servizi di calcolo serverless
    indipendenti e infrastruttura condivisa per database, data warehouse e servizi AI.}
  \label{fig:architettura_suite}
\end{figure}

% =====================================================================
\section{Vera --- assistente generale con Web grounding}
\label{sec:agenti-vera-arch}
% =====================================================================

\subsection{Caso d'uso e posizionamento}
\label{subsec:vera-caso-uso}

Vera è l'agente generalista del pacchetto, deve potere coprire i tipici casi d'uso per cui un utente potrebbe rivolgersi a un LLM esterno alla Cooperativa:
analisi e riassunto di documenti e testi, formazione e approfondimento, correzione di codice e formule e così via.
Deve quindi supportare la possibilità per l'utente di caricare documenti tramite tecnologia RAG~\cite{lewis2020rag} nel maggior numero possibile di formati e, per le richieste che necessitano di informazioni attuali, deve potere accedere
a internet attraverso il meccanismo di Web grounding~\cite{gcp_grounding} che ancora la generazione a risultati provenienti da Google Search e ne mitiga il rischio di allucinazione su fatti recenti. Quando le risposte fornite vengono tratte da documenti caricati o da internet, la fonte deve essere sempre presente e verificabile;
Ci deve essere certezza che nello svolgere queste ricerche non possa accadere esfiltrazione di dati presenti in documenti o prompt caricati dall'utente, che potrebbero essere confidenziali.

% =====================================================================
\section{CooPolicy --- RAG su corpus di policy condiviso}
\label{sec:agenti-coopolicy-arch}
% =====================================================================

\subsection{Caso d'uso e posizionamento}
\label{subsec:coopolicy-caso-uso}

CooPolicy è l'agente specializzato nella ricerca e nell'interpretazione di documenti normativi interni: comprende in sostanza tutta la documentazione presente sulla Intranet di Cooperativa,
che tratta gli argomenti più vari e si è espansa nel corso degli anni, fino a diventare frastagliata e spesso di non facile consultazione.
Le tematiche trattate dai documenti possono comprendere, a titolo di esempio: manuali d'uso, procedure, contratti di lavoro, organigrammi, comunicazioni organizzative...

A differenza di Vera, CooPolicy non deve avere accesso a fonti esterne: vincolo essenziale è che le risposte agli utenti siano fornite in maniera precisa e approfondita se il materiale è presente nel
contesto RAG fornito~\cite{lewis2020rag,guu2020realm}, oppure che non siano fornite affatto.
Anche in questo caso ogni fonte deve essere correttamente citata e accessibile.
Devono essere messe in atto tutte le tecniche più adeguate per instradare correttamente le richieste degli utenti anche se pervenute in formato molto stringato o interpretabile, attingendo allo stato dell'arte sulla espansione di query mediante LLM~\cite{ma2023query_expansion,gao2022hyde}.

% =====================================================================
\section{AllyCare / NPS --- Text-to-SQL sui dati dei sondaggi}
\label{sec:agenti-allycare-arch}
% =====================================================================

\subsection{Caso d'uso}
\label{subsec:allycare-caso-uso}

AllyCare è l'agente specializzato nell'analisi dei dati dei sondaggi NPS (Net Promoter Score) condotti dalla cooperativa.
La base dati di appoggio è in formato tabellare e sottoposta ad aggiornamento costante:
nei sondaggi,  oltre al chiedere un voto per ognuno dei vari aspetti dell'esperienza di acquisto, viene lasciata la possibilità di compilare
un campo di testo libero con ciò che il socio ritiene di segnalare.
Un flusso LLM dedicato analizza i commenti presenti e li cataloga in base a delle tabelle specifiche (reparto e argomento di riferimento, valutazione del sentiment del commento, e così via)
e infine compila una tabella "ai enriched" che comprende sia tutti i dati di base di ogni singola risposta (voti etc) che i nuovi "tag" generati.

L'agente AllyCare ha lo scopo di facilitare le analisi su questa tabella, svolgendo sia compiti che tradizionalmente avrebbero richiesto la creazione di reportistica ad hoc
(domande quantitive come "numero di voti positivi per il periodo X, per il negozio Y, e così via, oppure creazione di grafici di andamento nel tempo), sia compiti di analisi qualitativa
che richiederebbero molto tempo: domande come "elenca le maggiori criticità nella regione / negozio specifico", e così via.

La richiesta dell'utente deve quindi essere tradotta dal modello in una (o più) query SQL adeguate, che vengono poi eseguite e il cui risultato viene utilizzato
per fornire la risposta e le successive analisi. Il problema della traduzione di linguaggio naturale in SQL è oggetto di una linea di ricerca consolidata, con benchmark di riferimento come Spider~\cite{yu2018spider} e BIRD~\cite{li2023bird} e tecniche di decomposizione e auto-correzione che migliorano l'accuratezza su query complesse~\cite{pourreza2023dinsql}.




% =====================================================================
\section{SysAid --- Text-to-SQL sul ticketing in due varianti}
\label{sec:agenti-sysaid-arch}
% =====================================================================

\subsection{Caso d'uso}
\label{subsec:sysaid-caso-uso}

SysAid è l'agente specializzato nell'analisi della base dati di ticketing della Cooperativa;
Anche in questo caso, la base dati di appoggio è in formato tabellare: è una copia "in streaming" e in tempo reale del database effettivo del servizio.

L'agente ha lo scopo di facilitare le analisi su questa tabella, svolgendo in questo cas sia compiti che tradizionalmente avrebbero richiesto la creazione di reportistica ad hoc
(numeriche, andamenti, tempi di risoluzione...), sia compiti di analisi qualitativa che richiederebbero molto tempo: domande come "elenca le maggiori criticità nella regione / negozio specifico",
ma anche, data la presenza di dati realtime, possibilità di avere informazioni sullo stato di salute attuale dei sistemi informativi e inoltre di fornire aiuto al troubleshooting se impiegato
da personale del Service Desk: difatti l'agente può consultare velocemente tutta la base di ticket, recuperando con efficienza occorrenze precedenti di un dato problema e così via.

Anche in questo caso, la richiesta dell'utente deve quindi essere tradotta dal modello in una (o più) query SQL adeguate, che vengono poi eseguite e il cui risultato viene utilizzato
per fornire la risposta e le successive analisi.

\subsection{Diversificazione del vendor e confronto dei modelli}
\label{subsec:agenti-sysaid-confronto}

SysAid è l'unico agente della suite realizzato architetturalmente in due varianti parallele
e simultanee, alimentate da due fornitori di Large Language Model concorrenti (rispettivamente Google~\cite{gcp_vertexai} e Anthropic~\cite{anthropic_docs}).

Entrambe le varianti condividono la medesima base di conoscenza, la stessa
struttura di database, i medesimi strumenti per l'analisi dei dati e
l'interfaccia utente finale.

La motivazione principale di questa duplicazione architetturale è fornire un primo caso effettivo di implementazione agnostica di un agente interno, come da linee guida del progetto. La compresenza dei due modelli costituisce inoltre il presupposto sperimentale per una valutazione comparativa controllata della qualità delle risposte (cfr. Cap.~\ref{chap:orion-valutazione}), in linea con la metodologia LLM-as-a-Judge proposta in letteratura~\cite{zheng2023llmjudge}.

% =====================================================================
\section{Alternative scartate}
\label{sec:agenti-alternative}
% =====================================================================

Questa sezione documenta le principali alternative considerate e le
ragioni per cui non sono state adottate. Rendere espliciti i vicoli
ciechi esplorati contribuisce a posizionare le scelte effettuate nel
contesto delle opzioni disponibili.

\paragraph{Orchestratore centrale fin dalla versione~1.}
La prima alternativa valutata era introdurre un orchestratore
intelligente già nella versione iniziale della suite, sulla scorta dei
pattern multi-agente proposti dalla letteratura recente~\cite{wu2023autogen,guo2024multiagent,wang2024moa}, eliminando la
necessità per l'utente di selezionare esplicitamente l'agente. Questa
opzione è stata scartata per una ragione metodologica: costruire
contemporaneamente gli agenti specializzati e il loro orchestratore
avrebbe reso molto più difficile isolare e diagnosticare i problemi.
Validare ogni agente indipendentemente --- verificare che CooPolicy
recuperi i documenti giusti, che AllyCare generi SQL accurato, che
Vera citi correttamente le fonti --- è molto più semplice quando
l'agente è esposto direttamente all'utente. L'orchestratore è stato
quindi rimandato alla fase successiva, quando ogni agente aveva
raggiunto un livello di maturità sufficiente per essere integrato.
Questa scelta è la genesi di Orion.

\paragraph{Fine-tuning di un modello proprietario.}
Un'alternativa esplorata nella fase di analisi iniziale era addestrare
un modello sui dati interni della cooperativa, per creare un LLM
specializzato sul dominio aziendale. Questa opzione è stata scartata
per tre ragioni. In primo luogo, il costo di addestramento e
mantenimento di un modello su scala enterprise è proibitivo per
un'organizzazione che non ha nel proprio core business lo sviluppo
di modelli AI. In secondo luogo, il fine-tuning non risolve il
problema fondamentale del knowledge cutoff: un modello addestrato sui
dati di sei mesi fa non conosce i ticket aperti questa settimana, né
le policy aggiornate ieri. Il grounding --- RAG o Text-to-SQL --- è
la risposta corretta per i dati che cambiano nel tempo. In terzo
luogo, un modello proprietario richiede infrastruttura di inferenza
dedicata, che la cooperativa non possiede e il cui costo di
ownership non è giustificabile.

\paragraph{Deployment on-premise.}
L'alternativa di installare modelli open-source su infrastruttura
propria è stata valutata e scartata per ragioni di scala e costo.
Un deployment on-premise di modelli parametrici sufficientemente
capaci per uso enterprise richiederebbe un cluster GPU dedicato, con
i relativi costi di hardware, elettricità, raffreddamento e
manutenzione, che nessuna cooperativa con il profilo di
Coop Alleanza~3.0 può sostenere in modo economicamente razionale.
I servizi gestiti in cloud offrono un compromesso
qualità-costo-manutenzione largamente superiore per questo contesto
specifico.
